Fix \((\alpha,\beta,\gamma)\in\mathcal R_{\mathrm{low}}\).  The lemma
\(\mathrm{lem:low\mbox{-}region\mbox{-}exponent\mbox{-}ordering}\) gives
\[
0<p<t<q<\frac12.
\tag{1}
\]
Since \(t<q\), \(\mathrm{def:kbrw\mbox{-}rate}\) gives
\[
r_n^{\mathrm{KBRW}}=n^{-t}\vee n^{-q}=n^{-t}>0.
\tag{2}
\]
If the additional premise in the proposition held, division by the positive
quantity in (2) would give
\[
\frac{R_n^{\mathrm{band}}}{r_n^{\mathrm{KBRW}}}
\geq c n^{t-p}\longrightarrow\infty,
\]
where the limit follows from \(t-p>0\) in (1).

Without that additional premise,
\(\mathrm{thm:band\mbox{-}density\mbox{-}polynomial\mbox{-}frontier}\) yields only
\[
c n^{-q}\leq R_n^{\mathrm{band}}
\leq Cn^{-p}(\log n)^{2p}.
\]
After division by (2), this becomes
\[
c n^{t-q}leq
\frac{R_n^{\mathrm{band}}}{r_n^{\mathrm{KBRW}}}
\leq Cn^{t-p}(\log n)^{2p}.
\tag{3}
\]
The lower endpoint in (3) tends to zero and the upper endpoint tends to infinity.
More concretely, the three positive sequences \(n^{-q}\), \(n^{-t}\), and
\(n^{-p}\) all lie between the two rate endpoints for all sufficiently large
\(n\), while their ratios to \(n^{-t}\) tend respectively to zero, one, and
infinity.  Thus the available bracket alone determines no limit for the minimax
risk ratio.